\documentclass[10pt]{article}
\usepackage[margin=1in, paperwidth=8.5in, paperheight=11in]{geometry}
\usepackage{ifpdf,amsmath, amssymb, comment, color, graphicx, stmaryrd,setspace,enumitem,tikz, fancyhdr, wrapfig, textcomp, units, mathptmx, pdfpages, cancel, multicol}
\usepackage{hyperref}
\hypersetup{
    colorlinks=true,
    urlcolor=blue,
}

\setlength{\headheight}{14.5pt}
\newcommand{\Q}{\mathbb{Q}}
\newcommand{\R}{\mathbb{R}}
\newcommand{\Z}{\mathbb{Z}}
\newcommand{\vu}{\mathbf{u}}
\newcommand{\vv}{\mathbf{v}}
\newcommand{\vw}{\mathbf{w}}
\newcommand{\vi}{\mathbf{i}}
\newcommand{\vj}{\mathbf{j}}
\newcommand{\vk}{\mathbf{k}}
\newcommand{\vn}{\mathbf{n}}
\newcommand{\vr}{\mathbf{r}}
\newcommand{\va}{\mathbf{a}}
\newcommand{\vF}{\mathbf{F}}
\newcommand{\vL}{\mathbf{L}}
\newcommand{\vT}{\mathbf{T}}
\newcommand{\vN}{\mathbf{N}}
\newcommand{\proj}{\operatorname{proj}}
\newcommand{\orth}{\operatorname{orth}}
\newcommand\del\nabla
\newcommand\dotp[1][.5]{\,\mathbin{\vcenter{\hbox{\scalebox{#1}{$\bullet$}}}}\,}
\newcommand{\slice}{_\text{slice}}

% Solution text is in red. If you want the solutions to show, remove the \iffalse from the definition of the \red command.
%\newcommand{\red}[1]{ %\iffalse
%	\textcolor{red}{#1} }%\fi}

\newenvironment{red}{\color{red}}{\ignorespacesafterend}
\newenvironment{blue}{\color{blue}}{\ignorespacesafterend}
\newcommand{\green}[1]{\textcolor{green}{#1}}
\renewcommand{\section}[1]{\begin{center} \textbf{#1} \\\end{center}}
%
\hyphenpenalty=5000
\setlength{\parindent}{0in}
%\oddsidemargin=-.25in
\allowdisplaybreaks
\pagestyle{fancy}
\renewcommand{\headrulewidth}{0pt}
\lhead{MATH 202}
\rhead{Spring 2022}
%\lfoot{\copyright\ CLEAR Calculus 2010}
\cfoot{}

\begin{document}
%


%\onehalfspacing
\allowdisplaybreaks
%##################################################################
\section{PS\#11 - Series convergence - \red{Answer key} }

\everymath{\displaystyle}
\begin{enumerate}[leftmargin=0pt]

\item (\href{https://activecalculus.org/single/sec-8-3-series.html#CaV}{AC 8.3 Exercise 5}) In this exercise we investigate the sequence $\left\{\frac{b^n}{n!}\right\}$ for any constant $b$.
\begin{enumerate}
    \item Use the Ratio Test to determine if the series $\sum \frac{10^k}{k!}$ converges or diverges.
    
    \begin{red}
    To use the Ratio Test, we're going to look at the limit of the ratio of the next term to the current term $\lim_{n\to\infty} \frac{|a_{n+1}|}{|a_n|}$:
    \begin{align*}
        \lim_{n\to\infty} \frac{|a_{n+1}|}{|a_n|}
        &= \lim_{n\to\infty} \frac{10^{n+1}/(n+1)!}{10^n / n!} \\
        &= \lim_{n\to\infty} \frac{10^{n+1}}{(n+1)!}\cdot \frac{n!}{10^n} \\
        &= \lim_{n\to\infty} \frac{10 n!}{(n+1)!} 
         = \lim_{n\to\infty} \frac{10}{n+1} = 0.
    \end{align*}
    Because this limit is less than 1, the Ratio Test tells us that the series converges.
    \end{red}
    \item Now apply the Ratio Test to determine if the series $\sum \frac{b^k}{k!}$ converges for any constant $b$.
    
    \begin{red}
    There's nothing special about the number 10 in the work above; all our steps are just as valid if that's replaced with any other constant $b$:
    \begin{align*}
        \lim_{n\to\infty} \frac{|a_{n+1}|}{|a_n|}
        &= \lim_{n\to\infty} \frac{b^{n+1}/(n+1)!}{b^n / n!} \\
        &= \lim_{n\to\infty} \frac{b^{n+1}}{(n+1)!}\cdot \frac{n!}{b^n} \\
        &= \lim_{n\to\infty} \frac{b\cdot n!}{(n+1)!} 
         = \lim_{n\to\infty} \frac{b}{n+1} = 0.
    \end{align*}
    Because this limit is less than 1, the Ratio Test tells us that the series converges.
    \end{red}
    \item Use your result from (b) to decide whether the sequence $\left\{\frac{b^n}{n!}\right\}$ converges or diverges. If the sequence $\left\{\frac{b^n}{n!}\right\}$ converges, to what does it converge? Explain your reasoning.
    
    \begin{red}
    Since the series $\sum \frac{b^k}{k!}$ converges for any constant $b$, the Divergence Test tells us that the nth term has to go to 0. Therefore the sequence $\left\{\frac{b^n}{n!}\right\}$ must converge to 0.
    \end{red}
\end{enumerate}

\end{enumerate}
\end{document}